\documentclass[11pt]{article}
\usepackage[margin=1in]{geometry}
\usepackage{amsmath,amssymb,amsthm}
\usepackage{booktabs}
\usepackage{enumitem}
\usepackage{hyperref}
\hypersetup{colorlinks=true,linkcolor=blue,urlcolor=blue,citecolor=blue}

\theoremstyle{plain}
\newtheorem{theorem}{Theorem}[section]
\newtheorem{proposition}[theorem]{Proposition}
\newtheorem{corollary}[theorem]{Corollary}
\theoremstyle{definition}
\newtheorem{definition}[theorem]{Definition}
\newtheorem{example}[theorem]{Example}
\theoremstyle{remark}
\newtheorem*{remark}{Remark}

\title{2-d Kernels in Graphs and Digraphs: A Complete Walkthrough of FINDINGS.md}
\author{Companion explainer, written for a reader with no prior background\\ in kernel theory or in this project's tooling}
\date{August 21, 2026}

\begin{document}
\maketitle
\tableofcontents

\section{What this document is, and what it isn't}

FINDINGS.md is a lab notebook: dense, written for someone who already knows the
vocabulary, and organized around SQL queries and test counts rather than explanation.
This document is the opposite. Every symbol is defined before it is used, every
theorem gets a full proof written out in steps, and every experiment is explained in
plain language before its numbers are shown. Read this first; FINDINGS.md is the
underlying record and every number quoted here traces back to it.

Two phases of work sit behind this. Before any code existed, three facts about
\emph{digraphs} (Theorems~\ref{thm:A}, \ref{thm:B} and Example~\ref{ex:C} below) were
worked out by hand, on paper, from a counting argument and a parity argument. Then a
Python tool was built (Claude Code did the implementation) to test conjectures at
scale, and that tool both confirmed those hand-proofs on hundreds of thousands of
examples and produced four more theorems on its own (chordal graphs, simplicial
graphs, DAGs, and chordal-underlying digraphs), plus one refuted conjecture
(degenerate graphs) and one refuted conjecture that looked true up to 7 vertices and
broke at 8 (bipartite graphs). Section~\ref{sec:literature} also corrects a citation:
the general problem's hardness traces to a 1994 paper, not the 2015 paper this project
had been using.

The motivating context is a research direction toward improving list-coloring bounds;
that connection is being worked out separately and is not part of what's explained
here — this document only covers what has actually been proved or tested about
2-d kernels themselves.

\section{The basic vocabulary}

\subsection{Graphs, digraphs, edges and arcs}

A \textbf{graph} $G = (V, E)$ is a finite set of \textbf{vertices} $V$ together with a
set $E$ of \textbf{edges}, each an unordered pair of distinct vertices. Draw a vertex
as a dot and an edge as a line with no arrowhead: if $\{u,v\} \in E$ then $u$ and $v$
are \textbf{adjacent}, or \textbf{neighbours}. The number of vertices is the graph's
\textbf{order}, usually written $n$; the number of edges is its \textbf{size}, usually
written $m$. The \textbf{neighbourhood} of a vertex $v$, written $N(v)$, is the set of
vertices adjacent to it; its \textbf{degree} is $|N(v)|$.

A \textbf{digraph} (directed graph) $D = (V, A)$ replaces edges with \textbf{arcs}: an
arc is an \emph{ordered} pair $(u,v)$, drawn as a line with an arrowhead from $u$ to
$v$. Now $u \to v$ and $v \to u$ are different things, and a digraph may contain
neither, either, or both. The \textbf{out-neighbourhood} of $v$, $N^+(v)$, is the set
of vertices $v$ has an arc \emph{to}; the \textbf{in-neighbourhood} $N^-(v)$ is the set
$v$ has an arc \emph{from}. The \textbf{out-degree} $d^+(v) = |N^+(v)|$ and
\textbf{in-degree} $d^-(v) = |N^-(v)|$ are generally different numbers.
$\delta^+(D)$ denotes the \emph{smallest} out-degree over all vertices of $D$; the
condition ``$\delta^+ \ge 2$'' means every single vertex has at least 2 out-arcs.

A digraph has an \textbf{underlying graph}: forget the arrows and merge $u\to v$ or
$v \to u$ (or both) into a single undirected edge $\{u,v\}$. The underlying
neighbourhood of $v$, still written $N(v)$ when there's no ambiguity, is
$N^+(v) \cup N^-(v)$ — every vertex $v$ has an arc to or from, in either direction.
Note the containment $N^+(v) \subseteq N(v)$: this small fact does real work in
Section~\ref{sec:e8}.

\begin{definition}[Symmetric digraph and oriented graph]
Two special ways of turning arcs into edges, or edges into arcs, in opposite
directions:
\begin{itemize}[nosep]
\item A digraph is \textbf{symmetric} if whenever $u \to v$ is an arc, so is
$v \to u$. Every undirected graph $G$ can be turned into a symmetric digraph by
replacing each edge $\{u,v\}$ with \emph{both} arcs $u\to v$ and $v \to u$; this is how
this project treats every plain graph, so that one piece of code (built for digraphs)
handles both cases.
\item A pair of opposite arcs between the same two vertices, $u\to v$ together with
$v \to u$, is called a \textbf{digon}.
\item An \textbf{oriented graph} is a digraph with \emph{no} digons: for any two
vertices, at most one of $u \to v$, $v \to u$ is present. Equivalently, an oriented
graph is what you get by starting from an ordinary graph and choosing, for each edge,
a single direction to point it — the opposite extreme from a symmetric digraph, which
keeps both directions for every edge.
\end{itemize}
\end{definition}

Symmetric digraphs and oriented graphs sit at two ends of a spectrum: a symmetric
digraph has the most arcs possible for its underlying graph (both directions,
everywhere), an oriented graph has the fewest (exactly one direction, everywhere).
Section~\ref{sec:e3} is entirely about how much harder the oriented case is.

\subsection{Independent sets, domination, and the 2-d kernel}

\begin{definition}
A set $S \subseteq V$ is \textbf{independent} if no arc (in either direction) joins two
vertices of $S$. Equivalently, in the underlying graph, no edge has both endpoints in
$S$.
\end{definition}

An ordinary \textbf{kernel} of a digraph is an independent set $S$ such that every
vertex outside $S$ has \emph{at least one} arc into $S$ (\textbf{1-dominates}). This is
the classical notion — Richardson's theorem says any digraph with no odd directed
cycle has one — and it is mentioned here only for contrast.

\begin{definition}[2-d kernel — the object of this whole project]
A set $S \subseteq V$ is a \textbf{2-d kernel} of a digraph $D$ if:
\begin{enumerate}[nosep]
\item $S$ is independent, and
\item every $v \notin S$ satisfies $|N^+(v) \cap S| \ge 2$ — it has arcs to at least
two \emph{different} vertices of $S$ (\textbf{2-domination}).
\end{enumerate}
For an ordinary graph, treated as a symmetric digraph, this is exactly the
\emph{(2-d)-kernel} of A.~Wloch, \emph{On 2-dominating kernels in graphs}, Australas.
J. Combin.\ \textbf{53} (2012) 273--284 — the one piece of this problem that was
already named in the literature before this project started.
\end{definition}

A 2-d kernel is automatically a \textbf{maximal} independent set of the underlying
graph: nothing outside $S$ can be safely added to it, because everything outside $S$
is already adjacent to $S$ (it has at least 2 arcs into $S$, hence at least one edge
to $S$ in the underlying graph). This single observation is what makes a
brute-force reference solver possible at all (Section~\ref{sec:tool}).

\begin{example}[The smallest illustrative case: cycles]
\label{ex:cycle}
Let $C_n$ be the cycle on vertices $0, 1, \dots, n-1$, with edges between each
consecutive pair (and between $n-1$ and $0$) — every vertex has degree exactly 2.
Being 2-dominated by $S$ therefore means \emph{both} of a vertex's neighbours are in
$S$ (there's nothing else to draw a second arc from). So if $u \notin S$, both its
neighbours are in $S$ --- in particular, neither of $u$'s neighbours can also be
outside $S$, which means the vertices outside $S$ are themselves pairwise
non-adjacent. So \emph{both} $S$ and its complement are independent: together they
form a proper 2-colouring of the cycle. A cycle has a proper 2-colouring exactly when
$n$ is even (an odd cycle is the standard example of a graph that isn't bipartite),
and when it does, the colouring is unique up to swapping the two colours. So: $C_n$
has \textbf{no} 2-d kernel when $n$ is odd, and \textbf{exactly two} when $n$ is even ---
the two alternating classes of size $n/2$, and no others. This recurs constantly
below: the shortest odd cycle, $C_5$, turns out to be the single most common
obstruction found anywhere in this project.
\end{example}

\subsection{Cliques, chordal graphs, and simplicial vertices}

\begin{definition}
A \textbf{clique} is a set of vertices that are pairwise adjacent — every two of them
share an edge. An \textbf{induced subgraph} on a vertex subset $W$ keeps exactly the
vertices of $W$ and exactly the edges of the original graph that join two vertices of
$W$ — nothing added, nothing removed beyond what $W$ excludes.

A vertex $x$ is \textbf{simplicial} if its neighbourhood $N(x)$ is a clique — every two
neighbours of $x$ are themselves adjacent.

A graph is \textbf{chordal} if it has no induced cycle of length 4 or more (every
cycle of length $\ge 4$ has a ``chord'', an edge joining two non-consecutive vertices
of the cycle). The property that makes chordal graphs tractable throughout this
project is: \emph{every chordal graph with at least one vertex has a simplicial
vertex, and every induced subgraph of a chordal graph is itself chordal.} Both facts
are classical and are used, not re-derived, below.

A graph is a \textbf{simplicial graph} if every vertex lies in the closed neighbourhood
$N[x] = N(x) \cup \{x\}$ of some simplicial vertex $x$. This is a different, weaker
condition than chordal — it says every vertex is either simplicial itself or a
neighbour of one, without requiring the whole graph to be chordless-cycle-free.
\end{definition}

\subsection{Other graph classes that appear in the results}

These all appear as columns in the census (Section~\ref{sec:db}) and rows in the E1
and E2 tables (Section~\ref{sec:e1}). One sentence each:

\begin{center}
\begin{tabular}{@{}p{3.1cm}p{10.5cm}@{}}
\toprule
\textbf{Class} & \textbf{Plain-language definition} \\
\midrule
bipartite & vertices split into two groups with every edge crossing between them, never within a group \\
tree / forest & a tree is connected with no cycles at all; a forest is a disjoint union of trees \\
unicyclic & connected, with exactly one cycle \\
cactus & connected, every edge lies on at most one cycle \\
split & vertices split into a clique and an independent set (with edges allowed between them) \\
interval & vertices correspond to intervals on a line, edges exactly where two intervals overlap \\
cograph & no induced path on 4 vertices \\
claw-free & no induced ``claw'' (one central vertex joined to three mutually non-adjacent others) \\
planar & can be drawn in the plane with no edges crossing \\
regular & every vertex has the same degree \\
triangle-free & no three mutually adjacent vertices \\
block graph & every maximal 2-connected piece is a clique \\
DAG & \textbf{d}irected, \textbf{a}cyclic: a digraph with no directed cycle at all \\
\bottomrule
\end{tabular}
\end{center}

\subsection{DAGs, strong connectivity, and cycle parity}

\begin{definition}
A \textbf{directed cycle} is a sequence of arcs $v_1 \to v_2 \to \cdots \to v_k \to v_1$
using $k$ distinct vertices; its \textbf{length} is $k$. A digraph is \textbf{acyclic}
(a \textbf{DAG}) if it contains no directed cycle at all. Every DAG has a
\textbf{topological order}: a way of lining up the vertices left to right so every arc
points rightward. A \textbf{sink} is a vertex with no out-arcs — it sits at the far
right of every topological order.

A digraph is \textbf{strongly connected} (or just \textbf{strong}) if for every pair of
vertices $u, v$ there is a directed path from $u$ to $v$ \emph{and} one from $v$ to
$u$ — you can always get back to where you started. This is much stronger than being
\textbf{weakly connected}, which only asks that the underlying (undirected) graph be
connected.

``\textbf{All cycles even}'' describes a digraph in which every directed cycle has an
even number of arcs.
\end{definition}

\subsection{Degeneracy}

\begin{definition}
A graph is \textbf{$k$-degenerate} if every one of its subgraphs (not just induced
ones) has some vertex of degree at most $k$. Equivalently, you can remove the
vertices one at a time, always removing one of degree $\le k$ at the moment of its
removal, until nothing is left. The \textbf{degeneracy} of a graph is the smallest $k$
for which this works; it is computed by repeatedly peeling off a minimum-degree
vertex (the Matula--Beck algorithm). Degeneracy 0 means no edges at all; degeneracy 1
means the graph is a forest.
\end{definition}

\subsection{Isomorphism, canonical forms, and the graph6/digraph6 encoding}

Two graphs are \textbf{isomorphic} if one can be obtained from the other by
relabelling its vertices — they are ``the same shape'' even if drawn or listed
differently. An \textbf{isomorphism class} is the set of all labellings of one shape.

Telling two graphs apart up to isomorphism is itself a hard computational problem in
general. The tool \texttt{pynauty} (a Python wrapper around the \texttt{nauty}
library) computes a \textbf{canonical form} for any graph or digraph: a specific
labelling, chosen by a fixed rule, such that two graphs are isomorphic exactly when
their canonical forms are identical strings. This string is called a
\textbf{certificate}. Certificates are what let the census store one row per
\emph{shape}, not one row per \emph{labelling} — without them, the same graph drawn
with vertices numbered two different ways would be counted twice.

\textbf{graph6} and \textbf{digraph6} are short, fixed text encodings of a labelled
graph or digraph (a handful of ASCII characters per graph, growing slowly with $n$).
Combined with a canonical labelling, a graph6/digraph6 string becomes a compact,
human-typeable, unique name for an isomorphism class — this is the primary key used
in the database below.

\section{The tool: what was actually built}
\label{sec:tool}

\subsection{The verifier and the reference solver}

Every function in the codebase that produces a candidate set passes it through a
single verifier, \texttt{is\_2kernel}, before returning it: check independence, then
check that every outside vertex has at least 2 out-arcs into the set. This exists so
that no faster method can silently return a wrong answer — anything it hands back has
already been checked against the definition directly.

Because a 2-d kernel is always a maximal independent set of the underlying graph
(Section~2.2), a \textbf{complete, obviously-correct} solver is: enumerate every
maximal independent set (using the Bron--Kerbosch algorithm, a standard method for
this), and keep the ones that pass \texttt{is\_2kernel}. This is exponential in the
worst case, but it is the ground truth every faster method is checked against, and it
is cheap enough for everything in this project's exhaustive ranges.

\subsection{The forcing closure: five rules that sometimes decide the whole graph
for free}
\label{sec:forcing}

Before resorting to full search, a much cheaper propagation is tried first. It works
on a partial labelling of the vertices as \texttt{IN} (definitely in every 2-d kernel),
\texttt{OUT} (definitely in none), or \texttt{UNKNOWN} (not yet determined), applying
the following rules repeatedly until nothing changes:

\begin{description}[nosep,leftmargin=2.3cm,style=nextline]
\item[R0 (forced-in)] If $N^+(v)$ contains no independent pair (any two of $v$'s
out-neighbours are adjacent, or $v$ has fewer than 2 out-neighbours), then $v$ must be
\texttt{IN}. \emph{Why:} if $v$ were outside a 2-d kernel $S$, then $S$ would need two
non-adjacent members of $N^+(v)$ — impossible if no such pair exists.
\item[R1 (spread-out)] If $v$ is \texttt{IN}, every underlying neighbour of $v$ becomes
\texttt{OUT}. \emph{Why:} independence.
\item[R2 (conflict-check)] If $v$ is \texttt{OUT}, let $U(v)$ be the out-neighbours of
$v$ not yet marked \texttt{OUT}. If $U(v)$ contains no independent pair, declare
\texttt{CONFLICT} — no 2-d kernel can exist. \emph{Why:} $v$ can never find two
non-adjacent vertices to be 2-dominated by.
\item[R3 (forced-in by scarcity)] If $v$ is \texttt{OUT} and $|U(v)| \le 2$, every
member of $U(v)$ becomes \texttt{IN} (and if $|U(v)| < 2$, that is itself a
\texttt{CONFLICT}). \emph{Why:} $v$ needs 2 of them, and there are at most 2 left to
choose from.
\item[R4 (forced-in by the same test, on the fly)] If $v$ is \texttt{UNKNOWN} and
applying R2's test to $v$'s current non-\texttt{OUT} out-neighbours would give a
conflict, then $v$ itself must be \texttt{IN}. \emph{Why:} the same reasoning as R0,
recomputed as more vertices get decided.
\end{description}

Running these to a \textbf{fixed point} (repeat until no rule changes anything) gives
one of three verdicts: \texttt{CONFLICT} (proved no 2-d kernel exists), \texttt{SOLVED}
(every vertex decided, and the \texttt{IN} set is verified as an actual 2-d kernel), or
\texttt{UNDECIDED} (some vertices are left unresolved — this cheap method simply
didn't finish the job, and the full solver is needed).

\begin{example}[Forcing gets stuck]
\label{ex:stuck}
On the graph \texttt{G?KsZ\_} (Section~\ref{sec:e2}, the counterexample to the
bipartite conjecture): two pendant (degree-1) vertices, $0$ and $1$, fire R0
immediately (a degree-1 vertex's single neighbour trivially has no independent pair
among itself). That forces their neighbours $6$ and $7$ to \texttt{OUT} by R1. Now
$N(6)\setminus\text{OUT} = \{0,4,5\}$ and $N(7)\setminus\text{OUT} = \{1,2,3\}$; each
set has size 3 and contains a non-adjacent pair, so R2 doesn't fire (no conflict yet)
and R3 doesn't fire either (size is 3, not $\le 2$). Nothing else changes. The closure
stops with several vertices still \texttt{UNKNOWN} — an honest \texttt{UNDECIDED}, even
though (checked by the full solver) this graph in fact has no 2-d kernel at all.
\end{example}

Every disagreement of this kind found anywhere in the census (597 of them, across
almost 317{,}000 rows checked) is of exactly this shape: \texttt{UNDECIDED} when the
true answer is ``no 2-d kernel''. Not one case was found where the closure's
\texttt{CONFLICT} or \texttt{SOLVED} verdict was wrong — which makes sense, since both
of those verdicts come with their own certificate (a specific violated inequality, or
a verified set) baked into the rule that produced them.

\subsection{Recognizers and generators}

A separate module recognizes each of the classes in Section~2.4 (bipartite via
2-colouring, chordal via maximum cardinality search, and so on), and a generator
module produces the graphs and digraphs to test. Two kinds of generation are used:

\begin{itemize}[nosep]
\item \textbf{Exhaustive}: every graph up to a given size, with no gaps. Small cases
($n \le 7$) come from networkx's built-in atlas. Beyond that, \textbf{canonical
augmentation} is used: take every known graph on $n-1$ vertices, add one new vertex
joined to every possible subset of the existing ones, and keep only one representative
of each resulting isomorphism class (using pynauty's certificate to detect
duplicates). This is exhaustive because every graph on $n$ vertices, with any one
vertex deleted, becomes \emph{some} graph on $n-1$ vertices — so extending every
$(n-1)$-vertex graph in every possible way is guaranteed not to miss anything on $n$
vertices. This method was checked against known reference counts before being
trusted: it reproduces 12{,}346 graphs on 8 vertices, 274{,}668 on 9, and the standard
sequence of digraph counts (OEIS A000273) up to 1{,}540{,}944 digraphs on 6 vertices,
exactly.
\item \textbf{Random sampling}, with an explicit seed recorded for every run, used
only where exhaustive generation is too large (cubic graphs beyond small $n$, DAGs on
8 vertices).
\end{itemize}

\section{The database}
\label{sec:db}

\subsection{Why a plain file, not a server}

The results live in a single \textbf{sqlite} file. sqlite is not a server — there is no
process to start, no port, nothing a campus firewall could block — the ``database'' is
one file on disk, read and written directly by the Python standard library. This
matters in practice on a network that blocks locally-run servers (Postgres, MySQL):
sqlite simply has none to block.

\subsection{Schema}

Two tables. The first, \texttt{graphs}, has one row per distinct isomorphism class,
keyed by its canonical graph6/digraph6 string, with every computed property as a
column:

\begin{center}
\begin{tabular}{@{}p{3.6cm}p{10cm}@{}}
\toprule
\textbf{Column} & \textbf{What it holds} \\
\midrule
\texttt{key} (primary key) & the canonical graph6/digraph6 string \\
\texttt{n}, \texttt{m} & order and size \\
\texttt{degeneracy} & as defined in Section~2.6 (added and backfilled during E7) \\
\texttt{bipartite}, \texttt{chordal}, \texttt{split}, \texttt{interval}, \texttt{cograph}, \texttt{claw\_free}, \texttt{planar}, \texttt{regular}, \texttt{triangle\_free}, \texttt{block\_graph}, \texttt{simplicial\_graph}, \texttt{tree}, \texttt{forest}, \texttt{unicyclic}, \texttt{cactus}, \texttt{connected}, \texttt{strong}, \texttt{dag}, \texttt{symmetric}, \texttt{underlying\_bipartite}, \texttt{all\_cycles\_even} & one 0/1 flag column per class in Section~2.4, queried directly by name \\
\texttt{has\_2kernel} & 0/1, from the reference solver \\
\texttt{count\_2kernels} & how many distinct 2-d kernels it has \\
\texttt{min\_size}, \texttt{max\_size} & smallest and largest 2-d kernel, where any exist \\
\texttt{forcing\_verdict} & \texttt{CONFLICT} / \texttt{SOLVED} / \texttt{UNDECIDED} \\
\texttt{forcing\_agrees} & 0/1, whether that verdict matches \texttt{has\_2kernel} \\
\bottomrule
\end{tabular}
\end{center}

The second table, \texttt{membership}, is just pairs \texttt{(key, family)} — which
generation run (\texttt{atlas7}, \texttt{oriented6}, \texttt{cubic}, and so on,
Section~\ref{sec:e3}--\ref{sec:e7}) each graph came from. A graph can belong to more
than one family without its expensive properties (the 2-d kernel count, in particular)
ever being computed twice; every query in this document is a \texttt{JOIN} between the
two tables. Keying by canonical string also makes every run \textbf{idempotent} —
re-running a family after a bug fix updates the same row rather than creating a
duplicate.

\subsection{Reproducing everything}

\begin{verbatim}
python3.14 -m venv env && ./env/bin/pip install networkx pytest pynauty
./env/bin/python -m pytest -q
./env/bin/python -m twokernel.census run --family atlas7 --family graphs8 \
    --family graphs9 --family named --family oriented6 --family digraphs5 \
    --family digraphs6 --family dags7 --family dags8sample --family cubic \
    --family cubic_trianglefree --family cubic_girth5
./env/bin/python -m twokernel.census backfill degeneracy
\end{verbatim}
Every canned query in the sections below is printed in full, exactly as run, right
before its result.

\section{Testing philosophy}

Every claim below is one of exactly three things, and it is always labelled which:
\textbf{proved} (a full proof is given), \textbf{verified exhaustively} over a stated
finite range (every single instance in that range was checked by the reference
solver, not sampled), or an explicit \textbf{conjecture} with the checked range
stated. Known facts from the literature (Wloch's theorems, hand properties of the
Petersen graph) were encoded as regression tests; the rule throughout was that a
failing test is a bug to fix, never a value to edit.

\section{Result 1: how common are 2-d kernels? (FINDINGS.md \S2, ``E1'')}
\label{sec:e1}

Across \textbf{all 288{,}267 graphs on at most 9 vertices} (an exhaustive count, not a
sample), \textbf{125{,}439 (43.5\%)} have a 2-d kernel. The rate barely moves with $n$
(42.9\% at $n\le 7$ versus 43.5\% at $n \le 9$) but swings enormously by class:

\begin{center}
\begin{tabular}{@{}lrrl@{}}
\toprule
class & members & with a 2-d kernel & smallest members without one \\
\midrule
bipartite & 1{,}572 & 1{,}094 (70\%) & $K_2$ \\
chordal & 17{,}175 & 3{,}632 (21\%) & $K_2$ \\
split & 3{,}038 & 1{,}292 (43\%) & $K_2$ \\
interval & 12{,}657 & 2{,}511 (20\%) & $K_2$ \\
block graph & 1{,}575 & 211 (13\%) & $K_2$ \\
tree & 95 & 42 (44\%) & $K_2$, $P_4$ \\
planar & 87{,}835 & 30{,}343 (35\%) & $K_2$ \\
DAG & 10 & 10 (100\%) & none --- every DAG on $\le 9$ vertices in this range has one \\
\bottomrule
\end{tabular}
\end{center}

$K_2$ (two vertices, one edge) is the universal smallest obstruction: it belongs to
every class that allows an edge at all, and a 2-d kernel needs at least 2 members, which
$K_2$'s independence number of 1 can't supply. The only class in this table with
\emph{no} obstruction at all is DAGs — every one tested has a 2-d kernel, which
Section~\ref{sec:e6} explains and proves.

\section{Result 2: is the cheap method complete? (FINDINGS.md \S3, ``E2'')}
\label{sec:e2}

This is the central question: on which classes does the forcing closure of
Section~\ref{sec:forcing} decide every instance, with zero \texttt{UNDECIDED}
verdicts? Checked over all 288{,}267 graphs on at most 9 vertices, the classes that are
\textbf{always fully decided} are exactly: chordal, simplicial graphs, interval,
split, block graphs, forests, trees, and DAGs. Everything else — bipartite,
triangle-free, planar, cographs, cacti, and more — has at least one graph where
forcing stalls. The smallest failing example, in almost every class, is the same
graph: the 5-cycle $C_5$.

\subsection{Theorem: chordal graphs}

\begin{theorem}[E2.1]
\label{thm:e21}
On a chordal graph, the forcing closure always reaches \texttt{SOLVED} or
\texttt{CONFLICT} — never \texttt{UNDECIDED}. Consequently, a chordal graph has
\emph{at most one} 2-d kernel, and deciding 2-d kernel existence on chordal graphs is
solvable in polynomial time (the general problem, Theorem~\ref{thm:telle} below, is
NP-complete).
\end{theorem}

\begin{proof}
Run the closure to a fixed point, and suppose for contradiction some vertex is still
\texttt{UNKNOWN}. No \texttt{UNKNOWN} vertex can have an \texttt{IN} neighbour — R1
would already have marked it \texttt{OUT}. So for any \texttt{UNKNOWN} vertex $v$, the
set of its neighbours not yet marked \texttt{OUT} is exactly its neighbours that are
still \texttt{UNKNOWN} (since \texttt{IN} is impossible for a neighbour and everything
else is \texttt{OUT} or \texttt{UNKNOWN}). Let $W$ be the set of all \texttt{UNKNOWN}
vertices; it is non-empty by assumption. The subgraph induced on $W$ is an induced
subgraph of a chordal graph, hence itself chordal (Section~2.3), and every non-empty
chordal graph has a simplicial vertex — say $x \in W$ is one. Then $N(x) \cap W$ is a
clique (that is what simplicial means, restricted to $W$), so it contains no
independent pair, so R4's test fires on $x$: $x$ must be \texttt{IN}. But $x$ was
assumed \texttt{UNKNOWN} at the fixed point, where no rule can fire — contradiction.
So $W$ is empty: the closure decides everything.

With nothing \texttt{UNKNOWN}, the \texttt{IN} set at the end is the \emph{only}
candidate 2-d kernel the closure could possibly produce, so there is at most one.
\end{proof}

\subsection{Theorem: simplicial graphs}

\begin{theorem}[E2.2]
On a simplicial graph (Section~2.3), rules R0 and R1 alone decide the whole graph, and
again there is at most one 2-d kernel.
\end{theorem}

\begin{proof}
Every simplicial vertex has a clique neighbourhood, so R0 marks all of them
\texttt{IN} immediately. Every other vertex, by the definition of a simplicial graph,
lies in $N[x]$ for some simplicial $x$ — i.e., it is adjacent to an \texttt{IN} vertex
— so R1 marks it \texttt{OUT}. Nothing is left \texttt{UNKNOWN}.
\end{proof}

Theorem~\ref{thm:e21} accounts for chordal, split, interval, block graphs, forests,
and trees (all chordal); Theorem 3.2 covers simplicial graphs, which are \emph{not}
all chordal (attach one extra pendant vertex to every vertex of a 4-cycle: the result
is simplicial but not chordal, since the 4-cycle itself has no chord). These are
different theorems about different, overlapping classes. Both predict more than
completeness: they predict \emph{uniqueness}, and the census confirms it — the
maximum number of 2-d kernels found on either class is exactly 1, while cographs and
bipartite graphs reach as many as 4.

\subsection{Where this sits in the literature}
\label{sec:literature}

A 2-d kernel fits a general framework of Telle (1994): a $(\sigma,\rho)$-\emph{dominating
set} is an independent-flavoured set $S$ where members of $S$ have a number of
neighbours in $S$ from the set $\sigma$, and non-members have a number of neighbours
in $S$ from the set $\rho$. A 2-d kernel is exactly $\sigma = \{0\}$ (independence),
$\rho = \{2,3,4,\dots\}$ (2-domination).

\begin{theorem}[Telle, 1994]
\label{thm:telle}
Deciding whether a general graph has a $(\{0\}, \{q, q{+}1, q{+}2, \dots\})$-dominating
set is NP-complete, for every $q \ge 2$. Taking $q = 2$ shows 2-D KERNEL is NP-complete
on general graphs.
\end{theorem}

This is a citation, not a new proof — but it corrects one: this project had been
citing Bednarz, Hern\'andez-Cruz and Wloch, \emph{Ars Combin.}\ \textbf{121} (2015)
341--351 for this fact, which turns out to have been proved 21 years earlier, in J.~A.
Telle, \emph{Complexity of domination-type problems in graphs}, Nordic J. Comput.\
\textbf{1} (1994) 157--171, and the two lines of work do not appear to reference each
other.

A companion pair of papers by Golovach and Kratochv\'il proves a full dichotomy
(polynomial vs.\ NP-complete) of this same style of problem, restricted to chordal
graphs and separately to $k$-degenerate graphs. Their criterion has the same shape as
Theorem~\ref{thm:e21}: polynomial exactly when the graph class in question can never
have more than one such set. Whether their chordal dichotomy already covers the
2-d kernel's specific $\rho = \{2,3,\dots\}$ (an infinite, ``cofinite'' set, rather than
the finite $\rho$ their published abstract explicitly covers) could not be confirmed
--- both papers are paywalled and only their abstracts could be checked. Theorem~\ref{thm:e21}
is therefore presented here as independently found and proved, not claimed as a
corollary of a paper whose exact scope could not be verified.

\subsection{The bipartite conjecture, and why it is false}

At $n \le 7$, every one of 150 bipartite graphs was fully decided by forcing — as
clean a pattern as the chordal one. Extending the search to $n = 8$ broke it. The
counterexample, named \texttt{G?KsZ\_} in its canonical encoding, has 8 vertices, parts
$\{0,4,5,7\}$ and $\{1,2,3,6\}$, and two pendant vertices, $0$ and $1$, whose unique
neighbours ($6$ and $7$ respectively) lie in \emph{different} bipartition classes. It
has no 2-d kernel, and Example~\ref{ex:stuck} above traced exactly why forcing cannot
determine this on its own.

This is not a contradiction of Wloch's own even-distance theorem for bipartite graphs
(which needs every pair of pendant vertices to sit at even distance from each other) —
it is precisely the situation that theorem's hypothesis excludes, and it is exactly
the situation forcing cannot reconstruct locally. The lesson: two conjectures that
were indistinguishable on evidence up to 7 vertices turned out to be one theorem and
one false statement, and only extending the exhaustive range — not more theorizing —
told them apart.

\section{Result 3: digraphs (FINDINGS.md \S4, ``E3'')}
\label{sec:e3}

\subsection{The easy case: bipartite digraphs}

\begin{theorem}[A]
\label{thm:A}
If the underlying graph of a digraph $D$ is bipartite with parts $P$ and $Q$, and
every vertex of $Q$ has $d^+ \ge 2$, then $P$ is a 2-d kernel.
\end{theorem}

\begin{proof}
Bipartiteness means every arc runs between $P$ and $Q$, never within a part, so $P$ is
automatically independent. Every $v \in Q$ (everything outside $P$) has all of its
out-arcs landing somewhere, and since no arc can go to another vertex of $Q$, all of
$v$'s out-arcs land in $P$ — and there are at least 2 of them.
\end{proof}

This is nearly free (no digon restriction is even needed), and it sets up the harder,
non-bipartite generalization.

\subsection{Theorem B: strongly connected digraphs where every cycle is even}

\begin{theorem}[B]
\label{thm:B}
If $D$ is strongly connected, every directed cycle of $D$ has even length, and
$\delta^+(D) \ge 2$, then $D$ has two disjoint 2-d kernels.
\end{theorem}

\begin{proof}
Fix any vertex $r$. \emph{Claim: every directed path from $r$ to a given vertex $v$
has the same length parity (odd or even), regardless of which path is taken.} Suppose
$P$ and $Q$ are two $r \to v$ paths of different parity. Because $D$ is strongly
connected, there is some directed path $R$ from $v$ back to $r$. Both $P$ followed by
$R$, and $Q$ followed by $R$, are closed directed walks (they start and end at $r$).
Any closed directed walk in a digraph can be broken down into directed cycles whose
lengths add up to the walk's total length — this is a standard fact, obtained by
repeatedly cutting out a cycle wherever the walk revisits a vertex. Since every
directed cycle of $D$ has even length by hypothesis, both $|P| + |R|$ and $|Q| + |R|$
must be even (sums of even numbers). That forces $|P|$ and $|Q|$ to have the \emph{same}
parity (both equal to $-|R|$ modulo 2) — contradicting the assumption they differed.
So the claim holds: every vertex $v$ has a well-defined parity, $\mathrm{par}(v) \in
\{0,1\}$, equal to the length (mod 2) of any directed path from $r$ to it.

Let $V_0 = \{v : \mathrm{par}(v) = 0\}$ and $V_1 = \{v : \mathrm{par}(v) = 1\}$; these
partition $V$ since $D$ is strongly connected (every vertex is reachable from $r$).
\emph{Every arc flips parity}: if $u \to v$ is an arc, take any $r \to u$ path $P$ and
extend it by this arc to get an $r \to v$ path of length $|P| + 1$; by the claim,
\emph{every} $r \to v$ path (including this one) has the parity of $|P|+1$, so
$\mathrm{par}(v) = 1 - \mathrm{par}(u)$.

Two consequences. First, no arc can join two vertices of the same parity class (that
would contradict the flip just shown), so both $V_0$ and $V_1$ are independent.
Second, for $v \in V_i$, every out-neighbour of $v$ lies in $V_{1-i}$ (by the flip
rule) — so setting $S = V_{1-i}$, every vertex \emph{outside} $S$ (that is, every
vertex of $V_i$) has \emph{all} of its out-neighbours inside $S$, and there are at
least $\delta^+(D) \ge 2$ of them. So $V_{1-i}$ is a 2-d kernel, for each $i \in
\{0,1\}$ — giving two 2-d kernels, $V_0$ and $V_1$, which are disjoint since they
partition $V$.
\end{proof}

\subsection{Theorem C: strong connectivity cannot be dropped}
\label{sec:C}

\begin{example}[Theorem C's counterexample]
\label{ex:C}
Take vertices $\{1,2,3,4,a,b\}$. Put a symmetric 4-cycle on $1,2,3,4$ (arcs both ways
between each consecutive pair: $1\leftrightarrow2$, $2\leftrightarrow3$,
$3\leftrightarrow4$, $4\leftrightarrow1$), then add arcs $a\to1$, $a\to2$, $a\to b$,
$b\to2$, $b\to3$. Every directed cycle here is even: the digons ($1\leftrightarrow2$
etc.) are 2-cycles, and the outer 4-cycle $1\to2\to3\to4\to1$ has length 4; nothing
points \emph{into} $a$ or $b$, so neither can be part of any cycle. Every vertex has
out-degree at least 2 (check: each of $1,2,3,4$ has out-degree exactly 2 to its two
cycle-neighbours; $a$ has out-degree 3; $b$ has out-degree 2). But this digraph is
\textbf{not} strongly connected — nothing reaches $a$ or $b$ — and it has \textbf{no
2-d kernel at all}.

To see why: within $\{1,2,3,4\}$ alone (an even cycle, by Example~\ref{ex:cycle}) the
only two candidate 2-d kernels are $\{1,3\}$ and $\{2,4\}$.
Try $S \supseteq \{1,3\}$: vertex $a$ is then forced \texttt{OUT} (it has an arc to
$1 \in S$, and independence forbids $a$ joining $S$ alongside a vertex it points at),
so $a$ needs 2 out-neighbours in $S$ among $\{1,2,b\}$; it already has $1$, so it needs
$2$ or $b$ — but $2 \notin S$ in this branch, so $b$ must join $S$. But $b \to 3$ is an
arc and $3 \in S$, so $b$ joining $S$ violates independence. Contradiction. The
symmetric case, $S \supseteq \{2,4\}$, fails the same way ($a$ again forced
\texttt{OUT}, needing $b \in S$, but $b \to 2$ and $2 \in S$ blocks it). Both
possibilities fail, so no 2-d kernel exists.
\end{example}

Checked against the exhaustive digraph census: among all 126 digraphs on at most 6
vertices with $\delta^+ \ge 2$ that are ``all cycles even'' but \textbf{not} strongly
connected, exactly 84 (66.7\%) still have a 2-d kernel — so Example~\ref{ex:C} is a
typical failure mode when strong connectivity is dropped, not a contrived exception.
Restricting to the 81 digraphs in the same census that \emph{are} strongly connected
and all-cycles-even, \textbf{all 81} have a 2-d kernel, exactly as Theorem~\ref{thm:B}
guarantees.

\subsection{Theorem: the tight bound for oriented graphs}

Digons made Theorems~\ref{thm:A} and \ref{thm:B} easy — a kernel vertex could always
find its 2 required out-arcs, since nothing stopped it from pointing back at whoever
pointed at it. Forbidding digons (oriented graphs, Section~2.2) removes that freedom
entirely.

\begin{theorem}
No oriented graph on fewer than 8 vertices with $\delta^+ \ge 2$ has a 2-d kernel, and
$n = 8$ is attained.
\end{theorem}

\begin{proof}
Let $D$ be an oriented graph with $\delta^+ \ge 2$, and suppose it has a 2-d kernel $J$
with $|J| = j$; let $B = V \setminus J$, $|B| = t = n - j > 0$. Every vertex of $B$
sends at least 2 arcs into $J$ (2-domination), so summing over $J$: writing $d_x$ for
the number of $B$-vertices with an arc into $x \in J$,
\[
\sum_{x \in J} d_x \;\ge\; 2t.
\]
Now fix $x \in J$. Since $J$ is independent, every out-arc of $x$ goes to $B$. Since
$D$ is oriented (no digons), $x$ cannot have an out-arc to any vertex of $B$ that
\emph{already} points at $x$ — so $x$'s out-arcs are confined to the $t - d_x$
vertices of $B$ that do \emph{not} point at $x$. As $d^+(x) \ge 2$, this requires
$t - d_x \ge 2$, i.e.\ $d_x \le t - 2$ for every $x \in J$. Combining,
\[
2t \;\le\; \sum_{x\in J} d_x \;\le\; j (t-2).
\]
This forces $t \ge 3$ (else $t - 2 \le 0$ makes the bound impossible for a positive
$j$) and $j \ge \frac{2t}{t-2}$. Minimising $n = j + t$ over integers $t \ge 3$:

\begin{center}
\begin{tabular}{@{}rrrl@{}}
\toprule
$t$ & smallest $j$ satisfying $j \ge 2t/(t{-}2)$ & $n = j+t$ & \\
\midrule
3 & 6 & 9 & \\
\textbf{4} & \textbf{4} & \textbf{8} & minimum \\
5 & 4 & 9 & \\
6 & 3 & 9 & \\
7 & 3 & 10 & (and rising past this) \\
\bottomrule
\end{tabular}
\end{center}
So $n \ge 8$ always, with equality only possible at $t = j = 4$.
\end{proof}

\begin{example}[The bound is attained]
With $J = \{0,1,2,3\}$ and $B = \{4,5,6,7\}$, and arcs
\[
(0,5)(0,6)(1,6)(1,7)(2,4)(2,7)(3,4)(3,5) \quad\text{and}\quad
(4,0)(4,1)(5,1)(5,2)(6,2)(6,3)(7,0)(7,3),
\]
every vertex has out-degree exactly 2, no digon is present, and the 2-d kernels are
exactly $\{0,1,2,3\}$ and $\{4,5,6,7\}$.
\end{example}

\subsection{Exactly how many extremal examples are there? Two, not one}

The proof above pins $j = t = 4$ and forces \emph{equality} throughout: every kernel
vertex is pointed at by exactly $d_x = 2$ outside vertices, and (dually) every outside
vertex sends exactly 2 arcs into the kernel. That leaves one free choice per
outside-vertex — which 2 of the 4 kernel vertices it targets — subject to every kernel
vertex ending up targeted exactly twice overall; everything else is then forced (each
kernel vertex's own 2 out-arcs go to precisely the 2 outside vertices that do
\emph{not} target it). An exhaustive search over all $6^4 = 1296$ labelled choices
keeps exactly 90 that balance correctly, and canonical labelling splits those 90 into
exactly \textbf{2} isomorphism classes (of sizes 72 and 18). In both, the underlying
graph turns out to be the complete bipartite graph $K_{4,4}$ on $J \cup B$ — every
kernel vertex is adjacent (in the underlying sense) to \emph{all four} outside
vertices, not merely the two it directly exchanges arcs with — so both classes are
balanced orientations of $K_{4,4}$, differing only in \emph{which} orientation. They
are told apart by the pattern of target-pairs: the 72-example class (matching the
worked example above) uses the four ``consecutive'' pairs $\{0,1\},\{1,2\},\{2,3\},
\{3,0\}$ around a cyclic order of $J$; the 18-example class instead has the four
outside vertices split into two pairs that repeat a target-pair, e.g.\
$\{0,1\},\{0,1\},\{2,3\},\{2,3\}$.

\section{Result 4: cubic (3-regular) graphs and girth (FINDINGS.md \S5, ``E4'')}

A graph is \textbf{cubic} if every vertex has degree exactly 3. \textbf{Girth} is the
length of the shortest cycle in the graph; ``girth $\ge 5$'' rules out both triangles
and 4-cycles.

The question tested: does forbidding short cycles make a 2-d kernel more likely to
exist, among cubic graphs? Over seeded random samples at $n = 10, 12, \dots, 20$:

\begin{center}
\begin{tabular}{@{}lrrr@{}}
\toprule
family & isomorphism classes & with a 2-d kernel & fraction \\
\midrule
cubic, unrestricted & 293 & 160 & 0.546 \\
cubic, triangle-free & 240 & 131 & 0.546 \\
cubic, girth $\ge 5$ & 158 & 79 & 0.500 \\
\bottomrule
\end{tabular}
\end{center}

\textbf{Answer: no.} The fraction sits near one half regardless of girth, and does not
rise as the forbidden-cycle length grows. There is no support, at least among cubic
(degree-3) graphs, for the idea that avoiding short cycles helps guarantee a 2-d kernel.
This does not rule out the possibility for graphs of \emph{higher} regular degree — no
degree above 3 was tested here — but the degree-3 case gives that broader idea no
encouragement.

\section{Result 5: subdivisions (FINDINGS.md \S6, ``E5'')}

\begin{definition}
The \textbf{subdivision} $S(G)$ of a graph $G$ replaces every edge $\{u,v\}$ with a new
``subdivision'' vertex $w$ and two edges $\{u,w\}, \{w,v\}$ — every original edge
becomes a path of length 2 through a fresh vertex.
\end{definition}

\begin{theorem}
For every graph $G$, the original vertex set $V(G)$ is a 2-d kernel of its subdivision
$S(G)$.
\end{theorem}

\begin{proof}
In $S(G)$, no two original vertices are adjacent — every original edge has been
replaced by a path through a subdivision vertex — so $V(G)$ is independent. Every
other vertex of $S(G)$ is a subdivision vertex, and its neighbourhood is exactly the
two endpoints of the edge it replaced, both of which are in $V(G)$: exactly 2
neighbours in $V(G)$, satisfying 2-domination. (If $G$ has no edges at all, then
$S(G) = G$ and $V(G)$ is the whole vertex set, a 2-d kernel vacuously.)
\end{proof}

This gives a free, infinite family of graphs guaranteed to have a 2-d kernel — checked
directly on all 1253 graphs in the small-graph atlas (largest subdivision checked: 28
vertices) — and shows the property is never far away structurally: every graph is one
subdivision operation from having one.

\section{Result 6: DAGs (FINDINGS.md \S7, ``E6'')}
\label{sec:e6}

\begin{theorem}
A DAG has at most one 2-d kernel. Moreover, there is a linear-time procedure that either
outputs that unique 2-d kernel, or correctly certifies that none exists.
\end{theorem}

\begin{proof}
Fix a topological order $v_1, \dots, v_n$ of $D$ (every arc runs from a lower index to
a higher one), and process the vertices in reverse: $v_n, v_{n-1}, \dots, v_1$. When
$v_i$'s turn comes, every vertex of $N^+(v_i)$ has strictly higher index, so it has
already been permanently assigned a status, \texttt{IN} or \texttt{OUT} — nothing is
ever left ambiguous, which is the whole reason DAGs are tractable here.

Let $k = |N^+(v_i) \cap \{u : u \text{ already marked } \texttt{IN}\}|$ at the moment
$v_i$ is considered. Exactly one of three things happens, and it is forced with no
choice available:
\begin{itemize}[nosep]
\item \textbf{$k = 0$.} Being \texttt{OUT} needs $k \ge 2$, so $v_i$ cannot be
\texttt{OUT}. Being \texttt{IN} needs none of $v_i$'s out-neighbours to be
\texttt{IN}, which holds since $k=0$. Consistent: set $v_i := \texttt{IN}$.
\item \textbf{$k = 1$}, with the one \texttt{IN} out-neighbour called $u$. Being
\texttt{IN} is now impossible: the arc $v_i \to u$ together with $u$ already
\texttt{IN}, if $v_i$ joined it too, would put two adjacent vertices in the set,
violating independence. Being \texttt{OUT} is also impossible: that needs $k \ge 2$,
and $k=1$. \textbf{Neither status is available.} Since every 2-d kernel must give
$v_i$ one of the two statuses, and by the inductive hypothesis every 2-d kernel must
already agree with the statuses fixed on $N^+(v_i)$ so far, this means \emph{no
2-d kernel of $D$ exists at all}. Halt and report non-existence.
\item \textbf{$k \ge 2$.} Being \texttt{IN} is impossible: $v_i$ already has an
\texttt{IN} out-neighbour, so independence fails immediately. So $v_i$ must be
\texttt{OUT}, and that is consistent, since \texttt{OUT} needs exactly $k \ge 2$.
Set $v_i := \texttt{OUT}$.
\end{itemize}

The base case $v_n$ is automatically a sink (nothing can follow it in the order), so
$N^+(v_n) = \emptyset$ and $k=0$ trivially, giving $v_n := \texttt{IN}$ with no
contradiction; any other sink hits this same $k=0$ case whenever it is reached.

By strong induction from $i=n$ down to $i=1$, either the sweep completes with every
vertex assigned, or it halts at some $k=1$ vertex with a genuine proof of
non-existence, not a guess.

If it completes, let $S = \{v : v \text{ marked } \texttt{IN}\}$. Independence of $S$
now holds \emph{by construction}: if an arc $v_i \to v_j$ had both endpoints in $S$,
then when $v_i$ was set \texttt{IN} it must have been via the $k=0$ case, meaning
\emph{no} out-neighbour of $v_i$ — including $v_j$ — was \texttt{IN} at that time, so
$v_j$ was \texttt{OUT}, contradicting $v_j \in S$. And 2-domination of every
\texttt{OUT} vertex holds directly from the $k \ge 2$ condition that put it there,
which cannot change afterward since all of its out-neighbours were already
permanently decided. So $S$ is a genuine 2-d kernel, and since every step above was
forced with no branching, it is the only one.
\end{proof}

\begin{remark}
An earlier version of this proof merged the $k=0$ and $k=1$ cases into one (``fewer
than 2 \texttt{IN}, so set \texttt{IN}''), which is wrong exactly at $k=1$: setting
$v_i$ to \texttt{IN} there directly violates independence against its one already-\texttt{IN}
out-neighbour. The $k=1$ case is not a minor exception — it is precisely the
event whose growing probability, as $n$ increases, explains the falling fraction of
DAGs with a 2-d kernel reported below. There is no separate ``independence check''
tacked on at the end of the sweep; independence is guaranteed automatically whenever
the sweep completes, and failure is a specific, structural, locally-detected event
(some vertex landing at $k=1$), not a global property checked only after the fact.
\end{remark}


\begin{corollary}
Deciding whether a DAG has a 2-d kernel is not just polynomial, it is linear time —
answering, negatively, whether 2-D KERNEL might be NP-complete on DAGs the way finding
an ordinary kernel is (trivially) easy there too.
\end{corollary}

Confirmed on all 20{,}237 exhaustively-generated DAGs up to 7 vertices:
\texttt{MAX(count\_2kernels) = 1} throughout, and the forcing closure of
Section~\ref{sec:forcing} — which is not explicitly told about topological order at
all — independently decides every single one (0 \texttt{UNDECIDED}). One further
pattern in the data: the fraction of DAGs with a 2-d kernel falls steadily as $n$ grows
(100\%, 75\%, 50\%, 29.5\%, 16.9\% at $n = 1,4,5,6,7$) — by the corrected argument
above, this is the growing chance that \emph{some} vertex lands at $k=1$ (exactly one
already-\texttt{IN} out-neighbour) somewhere in the backward sweep, and nothing about
being acyclic protects against that as the graph grows.

\section{Result 7: $k$-degenerate graphs --- a refuted conjecture (FINDINGS.md \S8, ``E7'')}
\label{sec:e7}

The natural next question, given Theorem~\ref{thm:e21}'s chordal proof: does the same
argument survive \emph{weakening} chordal to low degeneracy (Section 2.6)? Every
chordal graph is, in particular, 1-degenerate at the low end, but degeneracy is a much
broader, weaker condition.

Over all 288{,}267 graphs on at most 9 vertices, grouped by exact degeneracy:

\begin{center}
\begin{tabular}{@{}rrrrr@{}}
\toprule
degeneracy & graphs & with a 2-d kernel & forcing disagreements & max \# 2-d kernels\\
\midrule
0 (edgeless) & 10 & 10 & 0 & 1 \\
1 (forests) & 299 & 109 & \textbf{0} & 1 \\
2 & 35{,}739 & 13{,}370 & \textbf{3{,}388 (9.5\%)} & 4 \\
3 & 173{,}728 & 74{,}357 & 34{,}839 & 4 \\
\bottomrule
\end{tabular}
\end{center}

\textbf{The conjecture is false, at the very next step past forests.} Forcing is
complete only at degeneracy $\le 1$ — and that case isn't new, since 1-degenerate
graphs are exactly the forests, which are already chordal and already covered by
Theorem~\ref{thm:e21}. As soon as degeneracy 2 is included, 9.5\% of graphs disagree,
and uniqueness fails too (up to 4 distinct 2-d kernels appear).

The smallest failing example is, once again, $C_5$ (the 5-cycle): every vertex has
degree exactly 2, so its degeneracy is 2 — but \emph{no} vertex of $C_5$ is simplicial,
since $C_5$ has no triangles at all, and a simplicial vertex needs its neighbourhood to
be a clique. This pinpoints exactly where the chordal proof's machinery breaks: it
relies on the existence of a \emph{simplicial} vertex (whose neighbourhood is a
clique) in every induced subgraph, which chordality guarantees but mere low degree
does not. A degree-2 vertex whose two neighbours are \emph{not} adjacent to each other
(as in $C_5$) is exactly the situation where R4 (Section~\ref{sec:forcing}) has
nothing to fire on. This is reported as a refuted conjecture, not an open one — an
explicit counterexample was found, not merely not-yet-found.

\section{Result 8: does the chordal theorem lift to digraphs? (FINDINGS.md \S9, ``E8'')}
\label{sec:e8}

Theorem~\ref{thm:e21} is about ordinary graphs. Does the same proof work for digraphs
whose \emph{underlying} graph is chordal?

\begin{theorem}
If the underlying graph of a digraph $D$ is chordal, then the forcing closure decides
$D$ completely, so $D$ has at most one 2-d kernel, and 2-D KERNEL is solvable in
polynomial time on such digraphs — with \emph{no} restriction on out-degree at all.
\end{theorem}

\begin{proof}
The one place the chordal proof (Theorem~\ref{thm:e21}) could plausibly break for
digraphs is R0/R4's test, which for a digraph looks at $N^+(v)$ (out-neighbours only),
whereas simpliciality is a property of the \emph{full} underlying neighbourhood
$N(v)$. But recall $N^+(v) \subseteq N(v)$ (Section~2.1): the out-neighbourhood is
always a \emph{subset} of the full neighbourhood. Since any subset of a clique is
itself a clique, if $N(v)$ is a clique (i.e.\ $v$ is simplicial) then $N^+(v)$ is
automatically a clique too — so every simplicial vertex fires R0/R4 regardless of
which arcs happen to point outward. (This shows simpliciality is, if anything, a
\emph{stronger} guarantee than the digraph rule needs, not a weaker one — the opposite
of what one might guess going in.)

With that step secured, Theorem~\ref{thm:e21}'s proof carries over unchanged: at a
fixed point with some vertex still \texttt{UNKNOWN}, the set $W$ of \texttt{UNKNOWN}
vertices induces a chordal subgraph (of the underlying graph), hence has a simplicial
vertex $x$ within $W$; $N(x) \cap W$ is a clique, so $N^+(x) \cap W \subseteq N(x)
\cap W$ is too, firing R4 on $x$ — contradicting the fixed point.
\end{proof}

Confirmed exhaustively over every digraph on at most 6 vertices with $\delta^+ \ge 2$
whose underlying graph is chordal: \textbf{251{,}991 digraphs, zero disagreements},
plus a broad random check (20{,}000 further digraphs, with no out-degree restriction at
all) finding the same zero disagreements — supporting the stronger, unrestricted
statement above.

\section{What this adds up to}

Four things stand out across all eight results.

First, the entire difficulty of this problem is concentrated in one structural
feature: the presence or absence of a simplicial vertex in every induced subgraph.
Forcing is not a heuristic that ``usually works'' — it is \emph{exactly} complete on
chordal graphs, simplicial graphs, and everything built from them, and it fails for
the first time on $C_5$, where no rule can fire on any vertex at all.

Second, plausible-looking patterns can be indistinguishable from theorems right up
until they aren't: the bipartite conjecture matched the chordal one perfectly at
$n \le 7$ and broke at exactly $n = 8$.

Third, forbidding digons (oriented graphs) is a far heavier restriction than it looks:
it turns $\delta^+ \ge 2$ from something that helps build a 2-d kernel (as in Theorems A
and B) into something that actively fights against it, forcing $n \ge 8$ before one can
even exist.

Fourth, the two attempted generalizations of the chordal theorem went in opposite
directions: weakening the hypothesis from chordal to merely low-degeneracy broke the
proof immediately (Result 7), while strengthening the object from a graph to a
digraph left the proof completely intact (Result 8) — and the reason in each case only
became clear by writing the actual proof argument out, not by guessing from the shape
of the analogy.

\section{Honest limitations}

\begin{itemize}[nosep]
\item Exhaustive ranges, stated plainly: all graphs to $n=9$; all digraphs with
$\delta^+\ge2$ to $n=6$; all oriented graphs with $\delta^+\ge2$ to $n=6$; all DAGs (of
the sink-or-$d^+\ge2$ kind) to $n=7$; all connected split graphs to $n=9$. Beyond
these, every number quoted is from seeded random sampling and is marked as such.
\item The cubic-graph results (Section 6) are samples, not an exhaustive count, though
each row does report distinct isomorphism classes rather than raw samples.
\item Counting the exact number of 2-d kernels (\texttt{count\_2kernels}) uses the
reference brute-force solver, so it is only reported where that solver is cheap
enough to run — the census's exhaustive ranges, not arbitrary large graphs.
\item Theorems in Sections~\ref{sec:e2}, \ref{sec:e6} and \ref{sec:e8} are fully
proved; the only conjecture in this whole document that could still have been
overturned by a wider search was the degeneracy one (Section 7), and it already was.
\item The citation check against Golovach and Kratochv\'il's dichotomy papers
(Section~\ref{sec:literature}) was done at the level of their published abstracts
only, since both papers are paywalled; the exact scope of their result (finite versus
cofinite $\rho$) could not be confirmed from primary text.
\end{itemize}

\section*{Quick notation reference}
\addcontentsline{toc}{section}{Quick notation reference}

\begin{center}
\begin{tabular}{@{}ll@{}}
\toprule
symbol & meaning \\
\midrule
$n$, $m$ & order (vertex count), size (edge/arc count) \\
$N(v)$ & neighbourhood of $v$ (underlying, undirected sense) \\
$N^+(v)$, $N^-(v)$ & out-neighbourhood, in-neighbourhood of $v$ in a digraph \\
$d^+(v)$, $d^-(v)$ & out-degree, in-degree of $v$ \\
$\delta^+(D)$ & the smallest out-degree over all vertices of $D$ \\
$S$ & a candidate or actual 2-d kernel \\
$K_n$ & the complete graph on $n$ vertices \\
$C_n$, $P_n$ & the cycle, the path, on $n$ vertices \\
\texttt{IN} / \texttt{OUT} / \texttt{UNKNOWN} & a vertex's status during forcing closure \\
\texttt{CONFLICT} / \texttt{SOLVED} / \texttt{UNDECIDED} & the closure's final verdict \\
\bottomrule
\end{tabular}
\end{center}

\end{document}
